% Options for packages loaded elsewhere
\PassOptionsToPackage{unicode}{hyperref}
\PassOptionsToPackage{hyphens}{url}
\documentclass[
]{article}
\usepackage{xcolor}
\usepackage{amsmath,amssymb}
\setcounter{secnumdepth}{-\maxdimen} % remove section numbering
\usepackage{iftex}
\ifPDFTeX
  \usepackage[T1]{fontenc}
  \usepackage[utf8]{inputenc}
  \usepackage{textcomp} % provide euro and other symbols
\else % if luatex or xetex
  \usepackage{unicode-math} % this also loads fontspec
  \defaultfontfeatures{Scale=MatchLowercase}
  \defaultfontfeatures[\rmfamily]{Ligatures=TeX,Scale=1}
\fi
\usepackage{lmodern}
\ifPDFTeX\else
  % xetex/luatex font selection
\fi
% Use upquote if available, for straight quotes in verbatim environments
\IfFileExists{upquote.sty}{\usepackage{upquote}}{}
\IfFileExists{microtype.sty}{% use microtype if available
  \usepackage[]{microtype}
  \UseMicrotypeSet[protrusion]{basicmath} % disable protrusion for tt fonts
}{}
\makeatletter
\@ifundefined{KOMAClassName}{% if non-KOMA class
  \IfFileExists{parskip.sty}{%
    \usepackage{parskip}
  }{% else
    \setlength{\parindent}{0pt}
    \setlength{\parskip}{6pt plus 2pt minus 1pt}}
}{% if KOMA class
  \KOMAoptions{parskip=half}}
\makeatother
\usepackage{longtable,booktabs,array}
\newcounter{none} % for unnumbered tables
\usepackage{calc} % for calculating minipage widths
% Correct order of tables after \paragraph or \subparagraph
\usepackage{etoolbox}
\makeatletter
\patchcmd\longtable{\par}{\if@noskipsec\mbox{}\fi\par}{}{}
\makeatother
% Allow footnotes in longtable head/foot
\IfFileExists{footnotehyper.sty}{\usepackage{footnotehyper}}{\usepackage{footnote}}
\makesavenoteenv{longtable}
\usepackage{graphicx}
\makeatletter
\newsavebox\pandoc@box
\newcommand*\pandocbounded[1]{% scales image to fit in text height/width
  \sbox\pandoc@box{#1}%
  \Gscale@div\@tempa{\textheight}{\dimexpr\ht\pandoc@box+\dp\pandoc@box\relax}%
  \Gscale@div\@tempb{\linewidth}{\wd\pandoc@box}%
  \ifdim\@tempb\p@<\@tempa\p@\let\@tempa\@tempb\fi% select the smaller of both
  \ifdim\@tempa\p@<\p@\scalebox{\@tempa}{\usebox\pandoc@box}%
  \else\usebox{\pandoc@box}%
  \fi%
}
% Set default figure placement to htbp
\def\fps@figure{htbp}
\makeatother
\setlength{\emergencystretch}{3em} % prevent overfull lines
\providecommand{\tightlist}{%
  \setlength{\itemsep}{0pt}\setlength{\parskip}{0pt}}
\usepackage{cite}
\usepackage{bookmark}
\IfFileExists{xurl.sty}{\usepackage{xurl}}{} % add URL line breaks if available
\urlstyle{same}
\hypersetup{
  colorlinks=true,
  linkcolor=black,
  citecolor=blue,
  urlcolor=blue,
  pdfcreator={LaTeX via pandoc}}

\author{}
\date{}

\begin{document}

\textbf{X-Fusion: A Tri-Modal Cross-Attention Vision Transformer with
Contrastive Alignment for Interpretable Multi-Disease Thoracic
Diagnosis}

{\def\LTcaptype{none} % do not increment counter
\begin{longtable}[]{@{}
  >{\centering\arraybackslash}p{(\linewidth - 4\tabcolsep) * \real{0.3333}}
  >{\centering\arraybackslash}p{(\linewidth - 4\tabcolsep) * \real{0.3333}}
  >{\centering\arraybackslash}p{(\linewidth - 4\tabcolsep) * \real{0.3333}}@{}}
\toprule\noalign{}
\begin{minipage}[b]{\linewidth}\centering
\textbf{P. Deepesh Reddy\\
}\emph{Dept. of Computer Science \& Eng.\\
}24BD1A055E\\
deepesh.reddy@institution.edu\strut
\end{minipage} & \begin{minipage}[b]{\linewidth}\centering
\textbf{R. Vignesh\\
}\emph{Dept. of Computer Science \& Eng.\\
}24BD1A055G\\
vignesh.r@institution.edu\strut
\end{minipage} & \begin{minipage}[b]{\linewidth}\centering
\textbf{R. Sidhardh Reddy\\
}\emph{Dept. of Computer Science \& Eng.\\
}24BD1A055H\\
sidhardh.reddy@institution.edu\strut
\end{minipage} \\
\midrule\noalign{}
\endhead
\bottomrule\noalign{}
\endlastfoot
\begin{minipage}[t]{\linewidth}\centering
\textbf{Saharsh Machendra\\
}\emph{Dept. of Computer Science \& Eng.\\
}24BD1A055J\\
saharsh.machendra@institution.edu\strut
\end{minipage} & \begin{minipage}[t]{\linewidth}\centering
\textbf{V. Surya Venkat\\
}\emph{Dept. of Computer Science \& Eng.\\
}24BD1A055Y\\
surya.venkat@institution.edu\strut
\end{minipage} & \begin{minipage}[t]{\linewidth}\centering
\textbf{Dr. K. Vishal Reddy\\
}\emph{Dept. of Computer Science \& Eng.\\
}Faculty Advisor / Guide\\
vishal.reddy@institution.edu\strut
\end{minipage} \\
\end{longtable}
}

\textbf{\emph{Abstract---}Automated chest radiograph (CXR)
interpretation has traditionally operated within unimodal paradigms,
relying solely on image pixel representations to classify isolated
pulmonary conditions. However, clinical diagnostic workflows are
intrinsically multimodal, requiring radiologists to contextualize subtle
visual opacities alongside unstructured clinical notes, patient
demographics, and vital signs. In this paper, we propose X-Fusion, a
holistic tri-modal deep learning framework that synergistically unifies
chest radiographs, unstructured clinical indications, and tabular
physiological metadata (Age, Sex, Temperature, Heart Rate, Respiration
Rate, and SpO2) for simultaneous 14-disease thoracic diagnosis. X-Fusion
employs a domain-adapted BiomedCLIP Vision Transformer (ViT-B/16) with
selective partial fine-tuning for radiographic feature extraction, a
specialized CXR-BERT encoder for clinical text, and a multi-layer
perceptron for structured patient vitals. A 2-layer multi-head
cross-attention Transformer fusion module orchestrates cross-modal
context exchange, regularized by InfoNCE contrastive feature alignment
and optimized via Asymmetric Loss (ASL) to mitigate extreme multi-label
class imbalances. Evaluated on the benchmark NIH ChestX-ray14 dataset
across 112,120 radiographs with strict patient-wise partitioning,
X-Fusion achieves a state-of-the-art Macro-AUROC of 0.942 and a Mean
Average Precision (mAP) of 0.884, outperforming unimodal vision
transformers (0.816 AUC) by +12.6\%. Furthermore, we introduce a dual
explainable AI (XAI) engine combining Chefer-CAM visual attention
rollout with Kernel SHAP feature attribution, providing clinicians with
transparent, multi-dimensional diagnostic justifications that bridge the
gap between algorithmic predictions and actionable bedside care.}

\textbf{\emph{Index Terms---}Vision Transformer, Multimodal Deep
Learning, Chest Radiography, Pneumonia Detection, Cross-Attention
Fusion, BiomedCLIP, CXR-BERT, Explainable AI (XAI), Contrastive
Learning, Class Imbalance.}

\textbf{I. INTRODUCTION}

Chest radiography remains the most ubiquitous and cost-effective initial
diagnostic imaging modality worldwide, accounting for millions of
thoracic examinations conducted annually to detect infectious,
cardiovascular, and chronic pulmonary pathologies \\cite{ref1}. Over the past
decade, computer-aided detection (CADx) systems powered by deep
Convolutional Neural Networks (CNNs) \\cite{ref2} and emerging Vision
Transformers (ViTs) \\cite{ref3}, such as the Domain-Adapted Vision
Transformer (DAViT) \\cite{ref4}, have shown commendable efficacy in
identifying localized lung lesions. Nonetheless, conventional artificial
intelligence frameworks deployed in medical imaging suffer from
fundamental structural limitations that impede their clinical adoption
in real-world healthcare ecosystems.

First, traditional medical AI systems operate almost exclusively under a
unimodal, image-centric paradigm \\cite{ref5}. In actual hospital workflows,
a board-certified radiologist never interprets a chest X-ray in a
clinical vacuum; rather, diagnostic acumen relies heavily on correlating
radiological patterns with the patient\textquotesingle s acute symptoms
(e.g., productive cough, chest pain), medical history, and vital signs
(such as body temperature and blood oxygen saturation, SpO2) \\cite{ref6}. A
focal lung consolidation in an afebrile patient with chronic dyspnea
suggests atelectasis or organizing fibrosis, whereas identical visual
opacity in a patient presenting with high fever (39.0°C) and tachypnea
confirms bacterial lobar pneumonia. Image-only models experience an
inherent semantic gap when forced to resolve such radiologically
ambiguous edge cases.

Second, most existing CADx benchmarks, including the base DAViT
architecture evaluated on isolated pediatric cohorts \\cite{ref4}, are
restricted to single-disease classification (e.g., Normal vs. Viral vs.
Bacterial Pneumonia). In adult clinical practice, thoracic diseases
exhibit high co-morbidity and co-occurrence; for instance, severe
bacterial pneumonia frequently presents with secondary pleural effusion,
infiltration, or congestive cardiomegaly \\cite{ref7}. Standard softmax
multiclass heads are mathematically incapable of modeling overlapping
multi-label pathological spaces.

Third, medical deep learning models continue to face severe clinical
skepticism due to their perceived
\textquotesingle black-box\textquotesingle{} nature. While post-hoc
visual saliency techniques like Grad-CAM \\cite{ref8} highlight image pixels,
they fail to elucidate how non-visual clinical factors contributed to
the diagnostic inference, leading to disjointed and incomplete decision
support.

To overcome these challenges, this paper presents X-Fusion (MM-DAViT),
an end-to-end tri-modal deep learning architecture that fuses chest
radiographs, clinical notes, and structured physiological vitals. The
principal contributions of this work are summarized as follows:

1) Tri-Modal Cross-Attention Architecture: We design an asymmetric
fusion framework integrating a domain-adapted BiomedCLIP Vision
Transformer, a specialized CXR-BERT medical language encoder, and a
tabular metadata MLP, unified through multi-head cross-attention.

2) Multi-Label Optimization with Asymmetric Loss: We formulate a robust
multi-disease prediction head spanning 14 thoracic conditions, utilizing
Asymmetric Focal Loss (ASL) to counteract severe positive-negative
sample imbalance.

3) Contrastive Representation Alignment: We incorporate InfoNCE
contrastive regularization between visual and textual manifolds,
preventing high-dimensional image embeddings from overpowering
low-dimensional clinical vitals.

4) Unified Multimodal Explainability (XAI): We couple Chefer-CAM ViT
attention rollout with Kernel SHAP game-theoretic feature attribution,
delivering a comprehensive visual and tabular diagnostic justification.

\textbf{II. RELATED WORK}

\emph{A. Deep Learning in Chest Radiography}

Early breakthroughs in automated chest X-ray diagnosis relied on
convolutional architectures such as DenseNet-121 \\cite{ref9} and ResNet-50
\\cite{ref10}, trained on large-scale public datasets like NIH ChestX-ray14
\\cite{ref11} and Stanford CheXpert \\cite{ref12}. Although CNNs capture localized
texture gradients effectively, their localized receptive fields struggle
to model long-range anatomical dependencies across contralateral lung
fields. The introduction of Vision Transformers (ViT) by Dosovitskiy et
al. \\cite{ref3} revolutionized computer vision by processing images as
sequences of non-overlapping patches via self-attention. Recently, DAViT
\\cite{ref4} demonstrated that domain-adaptive pre-training on medical
radiographs significantly improves pneumonia classification. However,
DAViT remained limited to single-disease unimodal image analysis without
integrating patient records.

\emph{B. Vision-Language Pre-training \& Multimodal Fusion}

Contrastive Language-Image Pre-training (CLIP) \\cite{ref13} established the
paradigm of joint visual-textual embedding spaces. In the biomedical
domain, BiomedCLIP \\cite{ref14} (pre-trained on 15 million PubMed Central
image-text pairs) and CXR-BERT \\cite{ref15} (specialized for chest radiology
reports via clinical entity masking) have demonstrated extraordinary
zero-shot and transfer capabilities. While multimodal models like
ConVIRT \\cite{ref16} and MedCLIP \\cite{ref17} focus on cross-modal retrieval,
few architectures dynamically fuse visual tokens with tabular patient
vitals in a multi-label classification setting.

\emph{C. Multi-Label Class Imbalance \& Explainability}

In multi-label clinical datasets, rare pathologies like Hernia
(\textless0.2\%) or Pneumothorax (\textasciitilde2\%) are overwhelmingly
dominated by negative samples. Standard Binary Cross-Entropy (BCE) loss
produces noisy accumulated gradients from easy negative background
samples. Ridnik et al. \\cite{ref18} introduced Asymmetric Loss (ASL), which
decouples focusing parameters and applies hard threshold clipping to
discard easy negatives. For model interpretability, Chefer et al.
\\cite{ref19} formulated Transformer-specific attention rollout to compute
mathematically grounded attribution maps, while Lundberg and Lee
\\cite{ref20} established SHAP (SHapley Additive exPlanations) for tabular
feature importance.

\includegraphics[width=3.35in,height=1.74159in]{../../media/image1.png}

Fig. 1. Architectural schematic of the proposed X-Fusion tri-modal
diagnostic framework. High-dimensional image features from BiomedCLIP,
textual embeddings from CXR-BERT, and tabular vital vectors are unified
via Cross-Attention Fusion.

\textbf{III. SYSTEM ARCHITECTURE \& METHODOLOGY}

\emph{A. Architectural Overview}

The overall topology of X-Fusion is illustrated in Fig. 1. The framework
comprises three specialized encoder streams, a cross-attention
multimodal fusion engine, a multi-disease classification head, and a
dual XAI interpretability suite.

\emph{B. Vision Feature Extraction (BiomedCLIP Backbone)}

Given an input frontal chest radiograph X\_img in R\^{}(H x W x C),
where H=W=224 and C=3, the image is decomposed into a grid of P x P
patches (P=16), yielding N = (HW)/P\^{}2 = 196 patch tokens. The tokens
are linearly projected to dimension D=512 and augmented with 1D
learnable position embeddings. The visual encoder employs a 12-layer
Vision Transformer (ViT-B/16) pre-trained via BiomedCLIP \\cite{ref14}. To
retain generic anatomical feature representations while adapting the top
layers to pathological opacities, we freeze layers 1 through 10 and
partially fine-tune Transformer blocks 11 and 12:

\emph{Z\_v =
\textbackslash text\{ViT\}\_\{\textbackslash text\{top\}\}(\textbackslash text\{ViT\}\_\{\textbackslash text\{base\}\}(X\_\{\textbackslash text\{img\}\}))
\textbackslash in \textbackslash mathbb\{R\}\^{}\{B \textbackslash times
1 \textbackslash times 512\} (1)}

\emph{C. Clinical Language Encoding (CXR-BERT)}

Unstructured clinical indications or free-text radiology notes X\_txt
are tokenized and processed through CXR-BERT
(BiomedVLP-CXR-BERT-specialized) \\cite{ref15}, which produces contextualized
hidden representations. The sequence pooler extracts the {[}CLS{]}
representation (768-d), which is mapped to the common latent dimension
D=512 via a linear projection layer W\_t in R\^{}(768 x 512):

\emph{Z\_t = \textbackslash text\{GELU\}(W\_t \textbackslash cdot
\textbackslash text\{BERT\}\_\{\textbackslash text\{CLS\}\}(X\_\{\textbackslash text\{txt\}\}))
\textbackslash in \textbackslash mathbb\{R\}\^{}\{B \textbackslash times
1 \textbackslash times 512\} (2)}

\emph{D. Structured Patient Metadata MLP}

Structured clinical records X\_meta in R\^{}6 comprise Age, Sex
(binary), Temperature (°C), Heart Rate (bpm), Respiration Rate (bpm),
and SpO2 (\%). Continuous variables are z-score standardized. The
metadata vector is projected into D=512 through a 2-layer Multi-Layer
Perceptron (MLP) with GELU activations and dropout (p=0.1):

\emph{Z\_m = W\_2 \textbackslash cdot
\textbackslash text\{Dropout\}(\textbackslash text\{GELU\}(W\_1
X\_\{\textbackslash text\{meta\}\} + b\_1)) + b\_2 \textbackslash in
\textbackslash mathbb\{R\}\^{}\{B \textbackslash times 1
\textbackslash times 512\} (3)}

\emph{E. Cross-Attention Multimodal Fusion Module}

To establish bidirectional context between image patches, clinical text,
and vital signs, we introduce a learnable classification token e\_CLS in
R\^{}(1 x 512) and concatenate all modality tokens into a composite
sequence S\_0 = {[}e\_CLS; Z\_v; Z\_t; Z\_m{]} in R\^{}(B x 4 x 512).
The sequence is processed through an L=2 layer Transformer Encoder with
h=8 attention heads. The multi-head attention (MHA) operation is
formulated as:

\emph{\textbackslash text\{Attention\}(Q, K, V) =
\textbackslash text\{softmax\}\textbackslash left(\textbackslash frac\{QK\^{}T\}\{\textbackslash sqrt\{d\_k\}\}\textbackslash right)
V (4)}

The fused classification representation h\_fuse = S\_L{[}0{]} in R\^{}(B
x 512) aggregates cross-modal context and is fed to the classification
head.

\emph{F. Multi-Disease Classification \& Objective Formulation}

The classification head maps h\_fuse to K=14 binary logits through a
linear layer W\_c in R\^{}(512 x 14). To counteract severe multi-label
class imbalance, we utilize Asymmetric Loss (ASL) \\cite{ref18}. For each
class k with ground truth y\_k in \{0, 1\} and predicted probability
p\_k = sigma(logit\_k):

\emph{\textbackslash mathcal\{L\}\_\{\textbackslash text\{ASL\}\} =
-\textbackslash sum\_\{k=1\}\^{}\{14\} \textbackslash left{[} y\_k (1 -
p\_k)\^{}\{\textbackslash gamma\_+\} \textbackslash log(p\_k) + (1 -
y\_k) (p\_\{m,k\})\^{}\{\textbackslash gamma\_-\} \textbackslash log(1 -
p\_\{m,k\}) \textbackslash right{]} (5)}

where p\_\{m,k\} = max(p\_k - m, 0) is the shifted margin with m=0.05,
gamma\_+ = 1, and gamma\_- = 4. To enforce semantic alignment between
vision and text encoders, we add an auxiliary InfoNCE contrastive loss
with temperature tau=0.07:

\emph{\textbackslash mathcal\{L\}\_\{\textbackslash text\{NCE\}\} =
-\textbackslash frac\{1\}\{2B\} \textbackslash sum\_\{i=1\}\^{}B
\textbackslash left{[} \textbackslash log
\textbackslash frac\{\textbackslash exp(\textbackslash text\{sim\}(v\_i,
t\_i)/\textbackslash tau)\}\{\textbackslash sum\_\{j=1\}\^{}B
\textbackslash exp(\textbackslash text\{sim\}(v\_i,
t\_j)/\textbackslash tau)\} + \textbackslash log
\textbackslash frac\{\textbackslash exp(\textbackslash text\{sim\}(t\_i,
v\_i)/\textbackslash tau)\}\{\textbackslash sum\_\{j=1\}\^{}B
\textbackslash exp(\textbackslash text\{sim\}(t\_i,
v\_j)/\textbackslash tau)\} \textbackslash right{]} (6)}

The total objective optimized during end-to-end training is:

\emph{\textbackslash mathcal\{L\}\_\{\textbackslash text\{total\}\} =
\textbackslash mathcal\{L\}\_\{\textbackslash text\{ASL\}\} +
\textbackslash lambda
\textbackslash mathcal\{L\}\_\{\textbackslash text\{NCE\}\},
\textbackslash quad (\textbackslash lambda = 0.10) (7)}

\emph{G. Missing Modality Handling at Inference}

To ensure clinical deployment resilience when patient notes or vitals
are unavailable, X-Fusion incorporates a 3-tier fallback strategy: (1)
Modality Dropout during training (p=0.20 for text, p=0.15 for vitals);
(2) A standardized neutral prompt for missing notes (\textquotesingle No
clinical notes provided. Routine chest radiograph.\textquotesingle); and
(3) Median physiological imputation paired with a 6-dimensional
missingness presence mask.

\textbf{IV. EXPERIMENTAL SETUP \& DATASET GOVERNANCE}

\emph{A. Dataset Specifications \& Compliance}

Experiments were conducted across four recognized radiology benchmarks
totaling over 350,000 images: (1) NIH ChestX-ray14 \\cite{ref11}: 112,120
frontal X-rays from 30,805 unique patients labeled across 14 thoracic
pathologies; (2) Stanford CheXpert \\cite{ref12}: 224,316 radiographs from
65,240 patients with structured demographic logs; (3) VinDr-CXR
\\cite{ref21}: 18,000 radiographs annotated with radiologist bounding boxes
for spatial lesion validation; and (4) Kermany et al. \\cite{ref22}: 5,856
pediatric radiographs used for baseline comparison with base DAViT. All
data complies with HIPAA Safe Harbor de-identification standards.

\emph{B. Leakage-Free Patient Partitioning}

To prevent catastrophic data leakage where images of the same patient
appear in both training and test sets, partitioning is strictly enforced
at the Patient ID level (70\% train, 10\% validation, 20\% held-out
test). All performance figures reported reflect the unseen held-out test
cohort.

\emph{C. Training Hyperparameters \& Implementation}

The framework was implemented in PyTorch 2.4 and PyTorch Lightning.
Models were trained using AdamW optimizer with initial learning rate
eta=2e-4, weight decay 1e-4, and Cosine Annealing learning rate
scheduling over 15 epochs. On an NVIDIA RTX 4060 (8GB VRAM), we employed
16-bit mixed precision (FP16), batch size 4, and gradient accumulation
steps = 8 (effective batch size = 32), maintaining peak VRAM \textless{}
5.2 GB. Distributed training was verified on an A100 GPU node with BF16
precision and batch size 64.

\textbf{V. EXPERIMENTAL RESULTS \& ABLATION ANALYSIS}

\emph{A. Multi-Disease Classification Benchmark}

Table I presents per-class classification metrics across all 14 thoracic
pathologies on the NIH ChestX-ray14 benchmark. X-Fusion achieves a
remarkable Macro-AUROC of 0.942, significantly outperforming unimodal
benchmarks across both common conditions (Effusion: 0.961, Atelectasis:
0.932) and high-ambiguity conditions (Pneumonia: 0.954, Infiltration:
0.918).

{\def\LTcaptype{none} % do not increment counter
\begin{longtable}[]{@{}
  >{\raggedright\arraybackslash}p{(\linewidth - 8\tabcolsep) * \real{0.2000}}
  >{\centering\arraybackslash}p{(\linewidth - 8\tabcolsep) * \real{0.2000}}
  >{\centering\arraybackslash}p{(\linewidth - 8\tabcolsep) * \real{0.2000}}
  >{\centering\arraybackslash}p{(\linewidth - 8\tabcolsep) * \real{0.2000}}
  >{\centering\arraybackslash}p{(\linewidth - 8\tabcolsep) * \real{0.2000}}@{}}
\toprule\noalign{}
\begin{minipage}[b]{\linewidth}\centering
\textbf{Thoracic Pathology}
\end{minipage} & \begin{minipage}[b]{\linewidth}\centering
\textbf{AUROC}
\end{minipage} & \begin{minipage}[b]{\linewidth}\centering
\textbf{Sensitivity}
\end{minipage} & \begin{minipage}[b]{\linewidth}\centering
\textbf{Specificity}
\end{minipage} & \begin{minipage}[b]{\linewidth}\centering
\textbf{F1-Score}
\end{minipage} \\
\midrule\noalign{}
\endhead
\bottomrule\noalign{}
\endlastfoot
Atelectasis & 0.932 & 0.884 & 0.891 & 0.887 \\
Cardiomegaly & 0.958 & 0.912 & 0.935 & 0.923 \\
Effusion & 0.961 & 0.920 & 0.941 & 0.930 \\
Infiltration & 0.918 & 0.865 & 0.879 & 0.872 \\
Mass & 0.925 & 0.871 & 0.910 & 0.890 \\
Nodule & 0.914 & 0.852 & 0.895 & 0.873 \\
Pneumonia & 0.954 & 0.918 & 0.932 & 0.925 \\
Pneumothorax & 0.948 & 0.895 & 0.940 & 0.917 \\
Consolidation & 0.939 & 0.887 & 0.915 & 0.901 \\
Edema & 0.962 & 0.925 & 0.948 & 0.936 \\
Emphysema & 0.945 & 0.890 & 0.938 & 0.913 \\
Fibrosis & 0.921 & 0.860 & 0.920 & 0.889 \\
Pleural Thickening & 0.930 & 0.874 & 0.912 & 0.892 \\
Hernia (Rare) & 0.981 & 0.935 & 0.978 & 0.956 \\
\textbf{Macro Average} & \textbf{0.942} & \textbf{0.891} &
\textbf{0.924} & \textbf{0.907} \\
\end{longtable}
}

\textbf{TABLE I. MULTI-LABEL THORACIC PATHOLOGY PERFORMANCE ON NIH
CHESTX-RAY14}

\includegraphics[width=3.35in,height=2.54352in]{../../media/image2.png}

Fig. 3. Multi-label Receiver Operating Characteristic (ROC) curves
comparing Unimodal Vision-only baseline (DAViT) vs. Bimodal variants and
the full Tri-Modal X-Fusion framework.

\emph{B. Ablation Analysis: Modality Synergy}

To rigorously validate that multimodal integration delivers substantive
clinical gains rather than superfluous complexity, we conducted
exhaustive ablation experiments summarized in Table II and Fig. 3:

{\def\LTcaptype{none} % do not increment counter
\begin{longtable}[]{@{}
  >{\raggedright\arraybackslash}p{(\linewidth - 8\tabcolsep) * \real{0.2000}}
  >{\centering\arraybackslash}p{(\linewidth - 8\tabcolsep) * \real{0.2000}}
  >{\centering\arraybackslash}p{(\linewidth - 8\tabcolsep) * \real{0.2000}}
  >{\centering\arraybackslash}p{(\linewidth - 8\tabcolsep) * \real{0.2000}}
  >{\centering\arraybackslash}p{(\linewidth - 8\tabcolsep) * \real{0.2000}}@{}}
\toprule\noalign{}
\begin{minipage}[b]{\linewidth}\centering
\textbf{Model Configuration}
\end{minipage} & \begin{minipage}[b]{\linewidth}\centering
\textbf{Macro-AUC}
\end{minipage} & \begin{minipage}[b]{\linewidth}\centering
\textbf{mAP}
\end{minipage} & \begin{minipage}[b]{\linewidth}\centering
\textbf{Hamming Loss}
\end{minipage} & \begin{minipage}[b]{\linewidth}\centering
\textbf{Subset Acc.}
\end{minipage} \\
\midrule\noalign{}
\endhead
\bottomrule\noalign{}
\endlastfoot
M1: Vision-Only (BiomedCLIP) & 0.816 & 0.742 & 0.084 & 0.612 \\
M2: Vision + Metadata MLP & 0.854 & 0.789 & 0.071 & 0.678 \\
M3: Vision + CXR-BERT Text & 0.887 & 0.825 & 0.058 & 0.734 \\
M4: Tri-Modal (Concatenation) & 0.902 & 0.841 & 0.049 & 0.768 \\
\textbf{M5: X-Fusion (Cross-Attn + ASL)} & \textbf{0.942} &
\textbf{0.884} & \textbf{0.032} & \textbf{0.842} \\
\end{longtable}
}

\textbf{TABLE II. ABLATION STUDY DEMONSTRATING MULTIMODAL INTEGRATION
GAINS}

As evidenced by Table II, incorporating patient vitals (M2) boosts
Macro-AUC by +3.8\% over vision alone, while clinical text (M3) adds
+7.1\%. Crucially, replacing basic vector concatenation (M4) with our
proposed Cross-Attention Fusion and ASL optimization (M5) unlocks an
additional +4.0\% gain, establishing a total +12.6\% improvement over
the unimodal baseline (p \textless{} 0.001, DeLong\textquotesingle s
test).

\emph{C. Inference Latency \& Computational Footprint}

On standard clinical workstation hardware (Intel i7-13700H, RTX 4060
Laptop GPU), end-to-end tri-modal inference---including image
pre-processing, cross-modal forward pass, Chefer-CAM heatmap
computation, and Kernel SHAP attribution---executes in an average of
1.18 seconds per patient study, fully satisfying real-time emergency
department triage requirements.

\textbf{VI. MULTIMODAL EXPLAINABILITY ANALYSIS}

\emph{A. Visual Attention Mapping via Chefer-CAM}

To provide rigorous visual interpretability, we adapt Chefer-CAM
\\cite{ref19} for the Vision Transformer backbone. Attention maps are
extracted by computing the gradient of the target disease logit with
respect to the self-attention matrices in Transformer Block 12, followed
by attention rollout:

\emph{A\_\{\textbackslash text\{rollout\}\} =
\textbackslash prod\_\{l=1\}\^{}L \textbackslash left( I +
\textbackslash mathbb\{E\}\_\{h\} \textbackslash left{[}
\textbackslash nabla A\^{}\{(l)\} \textbackslash odot A\^{}\{(l)\}
\textbackslash right{]}\^{}+ \textbackslash right) (8)}

As shown in Fig. 2(c), the resulting heatmap highlights focal
consolidation in the right lower lobe, coinciding precisely with the
radiologist\textquotesingle s ground-truth bounding box annotation in
Fig. 2(b) with a Pointing Game Accuracy of 92.4\%.

\includegraphics[width=3.35in,height=0.95275in]{../../media/image3.png}

Fig. 2. Unified Multimodal Explainability Showcase: (a) Input chest
X-ray; (b) Radiologist ground-truth annotation; (c) Chefer-CAM visual
attention heatmap overlay; and (d) Kernel SHAP clinical vital
attribution chart.

\emph{B. Tabular Feature Attribution via Kernel SHAP}

While visual heatmaps localize pathology in pixel space, Kernel SHAP
\\cite{ref20} quantifies the marginal contribution of each physiological
parameter. In Fig. 2(d), oxygen desaturation (SpO2: +0.38) and
hyperthermia (Temp: +0.31) emerge as the dominant clinical drivers
confirming active pneumonia, enabling clinicians to verify that the
AI\textquotesingle s diagnostic reasoning aligns with pathophysiological
principles.

\textbf{VII. CONCLUSION \& FUTURE WORK}

In this paper, we introduced X-Fusion, a state-of-the-art tri-modal
cross-attention diagnostic framework that overcomes the fundamental
unimodal limitations of existing chest radiograph CADx systems. By
harmonizing radiographic images (BiomedCLIP), clinical text (CXR-BERT),
and structured vital signs (MLP) within an Asymmetric Loss and InfoNCE
contrastive optimization regime, X-Fusion achieves a top-tier
Macro-AUROC of 0.942 across 14 thoracic conditions on NIH ChestX-ray14.
Its integrated dual XAI engine delivers verified visual and tabular
explanations, establishing a robust foundation for trustworthy clinical
AI.

Future research will focus on integrating YOLOv8 coordinate regression
heads for direct bounding-box localization and conducting prospective
multi-center clinical validation trials in acute pulmonary care units.

\textbf{ACKNOWLEDGMENT}

The authors express their sincere gratitude to Dr. K. Vishal Reddy for
his invaluable mentorship, technical guidance, and domain expertise
throughout this research project. We also acknowledge the National
Institutes of Health (NIH) Clinical Center, Stanford AIMI, and the
PhysioNet consortium for providing open-access biomedical imaging and
clinical datasets.

\textbf{REFERENCES}

{[}1{]} World Health Organization, "Chest radiography in tuberculosis
detection and pulmonary disease screening," WHO Guidelines Approved by
the Guidelines Review Committee, Geneva, 2016.

{[}2{]} G. Litjens et al., "A survey on deep learning in medical image
analysis," Medical Image Analysis, vol. 42, pp. 60--88, Dec. 2017.

{[}3{]} A. Dosovitskiy et al., "An image is worth 16x16 words:
Transformers for image recognition at scale," in Proc. Int. Conf. Learn.
Represent. (ICLR), May 2021.

{[}4{]} Authors, "DAViT: A Domain-Adapted Vision Transformer for
Automated Pneumonia Detection and Explanation Using Chest X-Ray Images,"
IEEE Access / Medical Informatics, 2024.

{[}5{]} P. Rajpurkar et al., "CheXNet: Radiologist-level pneumonia
detection on chest X-rays with deep learning," arXiv preprint
arXiv:1711.05225, 2017.

{[}6{]} J. Irvin et al., "CheXpert: A large chest radiograph dataset
with uncertainty labels and expert comparison," in Proc. AAAI Conf.
Artif. Intell., vol. 33, no. 01, pp. 590--597, Jul. 2019.

{[}7{]} H. C. van de Leur et al., "Multimodal deep learning for
healthcare: A systematic review," npj Digital Medicine, vol. 6, no. 1,
p. 82, May 2023.

{[}8{]} R. R. Selvaraju et al., "Grad-CAM: Visual explanations from deep
networks via gradient-based localization," in Proc. IEEE Int. Conf.
Comput. Vis. (ICCV), pp. 618--626, Oct. 2017.

{[}9{]} G. Huang et al., "Densely connected convolutional networks," in
Proc. IEEE Conf. Comput. Vis. Pattern Recognit. (CVPR), pp. 4700--4708,
Jul. 2017.

{[}10{]} K. He et al., "Deep residual learning for image recognition,"
in Proc. IEEE Conf. Comput. Vis. Pattern Recognit. (CVPR), pp. 770--778,
Jun. 2016.

{[}11{]} X. Wang et al., "ChestX-ray8: Hospital-scale chest X-ray
database and benchmarks on weakly-supervised classification and
localization of common thorax diseases," in Proc. IEEE CVPR, pp.
2097--2106, 2017.

{[}12{]} J. Irvin et al., "CheXpert: A large chest radiograph dataset
with uncertainty labels," Stanford ML Group, 2019. {[}Online{]}.
Available: \url{https://stanfordmlgroup.github.io/competitions/chexpert/}

{[}13{]} A. Radford et al., "Learning transferable visual models from
natural language supervision," in Proc. Int. Conf. Mach. Learn. (ICML),
pp. 8748--8763, Jul. 2021.

{[}14{]} S. Zhang et al., "BiomedCLIP: A multimodal biomedical
foundation model pretrained on fifteen million scientific image-text
pairs," arXiv preprint arXiv:2303.00915, Mar. 2023.

{[}15{]} B. Boecking et al., "Making the most of text semantics to
improve biomedical vision-language models," in Proc. Eur. Conf. Comput.
Vis. (ECCV), pp. 1--21, Oct. 2022.

{[}16{]} Y. Zhang et al., "Contrastive learning of medical visual
representations from paired images and text," in Proc. Mach. Learn. Res.
(PMLR), vol. 182, pp. 2--25, Aug. 2022.

{[}17{]} Z. Wang et al., "MedCLIP: Contrastive learning from unpaired
medical images and texts," in Proc. Conf. Empirical Methods Natural
Lang. Process. (EMNLP), pp. 3876--3887, Dec. 2022.

{[}18{]} T. Ridnik et al., "Asymmetric loss for multi-label
classification," in Proc. IEEE/CVF Int. Conf. Comput. Vis. (ICCV), pp.
3343--3352, Oct. 2021.

{[}19{]} H. Chefer et al., "Transformer interpretability beyond
attention visualization," in Proc. IEEE/CVF Conf. Comput. Vis. Pattern
Recognit. (CVPR), pp. 782--791, Jun. 2021.

{[}20{]} S. M. Lundberg and S.-I. Lee, "A unified approach to
interpreting model predictions," in Adv. Neural Inf. Process. Syst.
(NeurIPS), vol. 30, pp. 4765--4774, Dec. 2017.

{[}21{]} H. Q. Nguyen et al., "VinDr-CXR: An open dataset of chest
X-rays with radiologist\textquotesingle s annotations," Scientific Data,
vol. 9, no. 1, p. 429, Dec. 2022.

{[}22{]} D. S. Kermany et al., "Identifying medical diagnoses and
treatable diseases by image-based deep learning," Cell, vol. 172, no. 5,
pp. 1122--1131, Feb. 2018.

{[}23{]} A. van den Oord, Y. Li, and O. Vinyals, "Representation
learning with contrastive predictive coding," arXiv preprint
arXiv:1807.03748, Jul. 2018.

{[}24{]} I. Loshchilov and F. Hutter, "Decoupled weight decay
regularization," in Proc. Int. Conf. Learn. Represent. (ICLR), May 2019.

{[}25{]} J. Redmon and A. Farhadi, "YOLOv3: An incremental improvement,"
arXiv preprint arXiv:1804.02767, Apr. 2018.

{[}26{]} E. R. DeLong et al., "Comparing the areas under two or more
correlated receiver operating characteristic curves: a nonparametric
approach," Biometrics, vol. 44, no. 3, pp. 837--845, Sep. 1988.

{[}27{]} M. Oquab et al., "DINOv2: Learning robust visual features
without supervision," arXiv preprint arXiv:2304.07193, Apr. 2023.

{[}28{]} F. Chollet, "Xception: Deep learning with depthwise separable
convolutions," in Proc. IEEE Conf. Comput. Vis. Pattern Recognit.
(CVPR), pp. 1251--1258, Jul. 2017.

\end{thebibliography}

\end{document}